\documentclass[conference]{IEEEtran}
% \IEEEoverridecommandlockouts   % uncomment only if \thanks is needed
\usepackage{cite}
\usepackage{amsmath,amssymb,amsfonts}
\usepackage{algorithmic}
\usepackage{graphicx}
\usepackage{textcomp}
\usepackage{xcolor}
\usepackage{url}
\interdisplaylinepenalty=2500

\begin{document}

\title{Rate-Distortion Loss Attribution for AV1\\ Intra Partition Decisions}

\author{\IEEEauthorblockN{Pablo Rosa}
\IEEEauthorblockA{\textit{Video Technology Research Group (ViTech)} \\
\textit{Federal University of Pelotas} \\
Pelotas, Brazil \\
pablo.rosa@inf.ufpel.edu.br}
\and
\IEEEauthorblockN{Daniel Palomino}
\IEEEauthorblockA{\textit{Video Technology Research Group (ViTech)} \\
\textit{Federal University of Pelotas} \\
Pelotas, Brazil \\
dpalomino@inf.ufpel.edu.br}
\and
\IEEEauthorblockN{Marcelo Porto}
\IEEEauthorblockA{\textit{Video Technology Research Group (ViTech)} \\
\textit{Federal University of Pelotas} \\
Pelotas, Brazil \\
porto@inf.ufpel.edu.br}
\and
\IEEEauthorblockN{Luciano Agostini}
\IEEEauthorblockA{\textit{Video Technology Research Group (ViTech)} \\
\textit{Federal University of Pelotas} \\
Pelotas, Brazil \\
agostini@inf.ufpel.edu.br}
}

\maketitle

\begin{abstract}
The AOMedia Video 1 (AV1) codec decides how to split each intra-frame super block
through a recursive search that evaluates ten partition shapes per node, and
learned pruners are the dominant strategy to shorten that search. Every such
pruner couples the quality of the source that scores a node with the
aggressiveness of the policy acting on that score, so the reported time savings do
not identify which information carries the decision. This paper presents the
Representation Probe for Partitioning (RPP), an offline attribution protocol for
AV1 intra partitioning that reads the rate-distortion loss of competing
representations at matched search-cost reduction. RPP fixes one pruning policy,
sweeps its confidence threshold, and charges each forced decision the overhead
recorded by an exhaustive reference search, so that representations are compared
at the same amount of work skipped. Experimental results over 3.81 million
held-out decision nodes show that twenty-four compact descriptors reduce the loss
of isolated block variance by 7.0 times, that eight causal partitioning columns
reduce it by a further 1.60 times, and that four quantization and position columns
add nothing. A ConvNeXt with 28.1 million parameters reading raw pixels stays 4.1
times behind the compact descriptors, which hold 2481 times fewer parameters, and
doubling its fusion width makes it worse. The protocol is ordinal and does not
convert into Bjontegaard delta bitrate. To the best of the authors' knowledge,
this is the first work to attribute rate-distortion loss to information access in
AV1 partitioning under matched cost.
\end{abstract}

\begin{IEEEkeywords}
Video coding, AV1, intra-frame partitioning, feature representation, model
capacity, machine learning.
\end{IEEEkeywords}

\section{Introduction}

Video traffic keeps its position as the dominant share of Internet usage, and the
proportion of users consuming online video every month has remained above ninety
percent since 2022 \cite{statista}. Two families of codecs answer that demand: the
standardization path produced High Efficiency Video Coding \cite{hevc} and
Versatile Video Coding \cite{vvc}, both developed jointly by ISO/IEC MPEG and
ITU-T VCEG, while the Alliance for Open Media released the AOMedia Video 1 (AV1)
in 2018 as an open and royalty-free codec \cite{av1site}, drawing on VP9, Thor,
and Daala.

AV1 follows the hybrid block-based coding flow, combining intra-frame and
inter-frame prediction, transforms, quantization, in-loop filters, and entropy
coding, and it introduced new tools at every one of those stages
\cite{han2021,chen2018}. Among them, the partition structure was extended to ten
shapes per node, against the four of VP9. That coding efficiency is paid for in
computational effort, and comparative studies report AV1 encoding times far above
those of its competitors \cite{grois2016,layek2017}, with the gap widening at the
highest quality configuration.

Inside intra-frame coding, the recursive partition search is the single most
expensive stage at that configuration. The search visits a quad-tree of square
nodes, evaluates up to ten shapes at each one through full rate-distortion
optimization \cite{sullivan1998}, and descends into every node it does not
resolve, so the cost grows with the number of nodes rather than with the
difficulty of the content.

Several works attack this cost with learned decisions. In AV1, tool selection
guided by user priority and by coding efficiency was proposed in
\cite{xu2020,xu2021}, a fast intra mode decision based on the accuracy of a
rate-distortion model in \cite{jeong2019}, and decision trees for inter prediction
in \cite{kim2019}. The same problem in Versatile Video Coding was addressed with
light gradient boosting \cite{saldanha2022}, random forests \cite{he2021}, and
variance and gradient heuristics \cite{chen2019}, while dedicated hardware
architectures targeted the AV1 intra prediction modes \cite{correa2020,netto2022}.
These works can be classified as coding-efficiency oriented
\cite{jeong2019,xu2021,saldanha2022,he2021,chen2019}, hardware oriented
\cite{correa2020,netto2022}, or aimed at a different coding stage
\cite{xu2020,kim2019}. All of them report time savings obtained under a policy of
their own choosing, and none of these works separates the quality of the scoring
source from the aggressiveness of the policy, nor measures which information is
decisive for the partition choice.

This paper presents the Representation Probe for Partitioning (RPP), an offline
attribution protocol for AV1 intra partitioning. The main strategy of this
solution is to hold the pruning policy and the amount of skipped search work
fixed, and to vary only what a representation is allowed to see, so that the
resulting difference in rate-distortion loss is attributable to information access
alone. Under this protocol, compact descriptors outperform a high-capacity
convolutional model reading raw pixels, and the causal partitioning neighborhood
carries decision information that luminance does not provide.

\section{Intra Partition Search in AV1}

AV1 encodes each frame as a grid of super blocks of 64$\times$64 samples, and each
super block is recursively divided by a search over a quad-tree of square nodes of
64, 32, 16, and 8 samples. At every node the encoder may evaluate ten partition
shapes, organized in four groups: the undivided block, the quaternary split, the
two rectangular shapes, and the six extended shapes.

The evaluation order inside a node is fixed. The undivided candidate is evaluated
first and produces a complete rate-distortion cost; the quaternary split follows,
and opening it triggers the recursion into the four children; the rectangular
shapes come next, and the extended shapes last.

This order defines when each kind of information becomes available, which is the
property this work exploits. Before the search of a node begins, the encoder
already holds the source luminance of the block, the frame quantization index, the
position of the node, and the partition decisions of the causal neighbors above
and to the left, but it does not yet hold any rate-distortion cost for the current
node. Fig.~\ref{fig_timeline} summarizes this timeline.

The action studied in this paper is the earliest one available: committing the
node to the undivided shape before any rate-distortion evaluation is paid, and
therefore skipping every remaining candidate of the node and the whole subtree
below it. This single action has two measurable consequences, the search work that
is no longer performed and the rate-distortion cost that is lost whenever the
discarded subtree was the optimal choice.

This work is focused on that commitment at the pre-search insertion point of the
intra-frame partition tree. Inter-frame prediction, transform type selection,
quantization, in-loop filters, and entropy coding are not included in the loop
attacked here, and therefore they are not focused on this work. Blocks of
8$\times$8 samples are terminal leaves whose reference label is invariably the
undivided shape, and they are excluded from modeling while kept in the cost
accounting.

\begin{figure}[!t]
\centering
% TODO: replace by fig1_information_timeline.pdf, generated last in av1_bench.
\framebox[0.98\columnwidth]{\rule{0pt}{3.6cm}%
\parbox{0.86\columnwidth}{\centering\footnotesize
PLACEHOLDER --- information-availability timeline inside one node: causal context
and source luminance available before the search; undivided cost available after
the first candidate; subtree costs available only after the recursion. The
pre-search instant is marked.}}
\caption{Information available to a pruner at each instant of the search of a
partition node, and the pre-search instant addressed in this work.}
\label{fig_timeline}
\end{figure}

\section{Proposed RPP Protocol}

RPP is a measurement protocol, not an encoder solution. It compares
representations by what they are allowed to see, under one fixed policy and at a
matched amount of skipped search work.

\subsection{Representations Declared by Information Access}

The representations are declared by information access rather than by
architecture. The first is the isolated variance of the block, the
classical descriptor of partition heuristics. The second, denoted A, is a
twenty-four column vector of compact luminance descriptors of the block and of its
hierarchical context: variance and its per-quadrant dispersion, gradient energy
and orientation, row and column profiles, edge statistics, direct-current level,
and the contrast against the parent and sibling blocks. Three of those columns do
not derive from luminance, namely the normalized quantization index and the two
coordinates of the node inside the 64-sample unit.

Block B adds eight columns of causal partitioning neighborhood: availability of
the above and left neighbors, their already decided widths and heights in
logarithmic scale, the relative granularity of the neighborhood, and its
anisotropy. Block C adds four columns of effective quantization and position,
namely the dequantization step of the direct-current coefficient, the normalized
position within the frame, and the depth of the node. Four rungs are built from
these blocks, namely A, A+B, A+C, and A+B+C, the last one holding thirty-six
columns. All of these columns are natively available during encoding at the
pre-search instant, ensuring zero computational overhead for feature extraction.

The remaining representation is the deep pixel arm, a multi-level convolutional
model of the ConvNeXt family \cite{convnext} that mirrors the partition hierarchy.
Its input is a two-channel tensor of 64$\times$64 samples, formed by the
super-block luminance and a constant quantization plane. Three stages of the trunk
are merged top-down, so that each cell of the resulting map holds both its
immediate neighborhood and the context of the whole super block, and per-level
heads read average poolings of it. This arm is present as a test of the hypothesis
that raw pixels plus representation capacity suffice, and not as an architectural
baseline. A random scorer is included as a control.

The four rungs are multilayer perceptrons with two hidden layers of 64 and 32
rectified units, one network per block size, trained with hard-label cross entropy
and AdamW \cite{adamw} at a learning rate of $1 \times 10^{-3}$, batches of 4096
samples, and a fixed budget of 30 epochs. Every rung is trained with this one
recipe and differs only in which columns it reads, which is what makes the verdict
fall on information rather than on capacity or on training procedure. Each rung is
trained with three seeds and the reported value is their mean.

\subsection{Policy, Cost Model, and Loss}

The policy is fixed for every representation. A node $n$ is committed to the
undivided shape whenever the score assigned to that class exceeds a confidence
threshold $\tau$, in which case the node contributes only the cost of its
undivided candidate and the recursion stops; otherwise the node is evaluated in
full and the search descends. The rate-distortion loss charged to a committed node
is the overhead of that commitment against the optimal subtree, as recorded by the
exhaustive reference search,
\begin{equation}
r(n) = J_{\mathrm{none}}(n) - J^{*}(n),
\label{eqn_regret}
\end{equation}
where $J_{\mathrm{none}}(n)$ is the rate-distortion cost of the undivided
candidate and $J^{*}(n)$ is the cost of the optimal subtree rooted at $n$. By
construction $r(n) \geq 0$, and it is zero whenever the undivided shape already
was the reference decision. No encoding is repeated, since the whole protocol is a
counterfactual replay over the log of one exhaustive search.

The search work of a node of dimension $n$ is counted analytically as
\begin{equation}
c(n) = k_{n}\,n^{2}, \quad k_{n} = 9 \ \mathrm{for} \ n \in \{64,32,16\},
\ k_{8} = 4,
\label{eqn_cost}
\end{equation}
that is, the number of shape candidates a full search would evaluate, weighted by
the area of the block. The cost reduction $\Delta(\tau)$ is the percentage
difference between the work performed by the policy and the work of the full
search. It is a counted quantity and not wall-clock time, and it is not an
estimate of wall-clock time either.

The rate-distortion loss of a representation at a threshold $\tau$ is
\begin{equation}
L(\tau) = 100 \, \frac{\sum_{n \in P(\tau)} r(n)}{\sum_{n \in N}
J_{\mathrm{none}}(n)},
\label{eqn_loss}
\end{equation}
where $P(\tau)$ is the set of committed nodes and $N$ is the set of all decision
nodes. The denominator is accumulated before any threshold is fixed and stays
constant along the sweep, which is the condition for $L$ to be readable as a
frontier, since it tends to zero as pruning tends to zero.

Two properties of $L$ restrict its reading, and they are stated before any number
is reported. First, the denominator aggregates the undivided costs of the three
hierarchy levels, which cover the same image area, so it is systematically larger
than the cost of coding the frame and $L$ is systematically smaller than any
efficiency loss observable in the encoder. Second, the replay discounts recorded
values, whereas pruning inside the encoder also alters the reconstruction that
serves as reference to later blocks, a propagation absent here. Consequently, no
conversion factor between $L$ and the Bj\o ntegaard delta bitrate
\cite{bjontegaard} is postulated, and $L$ is used strictly in an ordinal regime,
in which every ratio is taken between representations read at the same operating
point. Each representation therefore yields a curve rather than a point, and the
reference reading in this paper is $\Delta = 25\%$.

\subsection{Dataset and Experimental Setup}

The reference labels were extracted from the libaom v3.10.0 reference software
\cite{libaom}, instrumented under a compile-time guard so that the production
binary remains byte-identical to the original. The encoder was configured for
All-Intra following the AOM Common Test Conditions \cite{ctc} at
\texttt{cpu-used=0}, which enables the full rate-distortion search and therefore
makes the recorded decisions optimal by construction.

The corpus comprises sixteen UVG sequences \cite{uvg} at 4K resolution
(3840$\times$2160), 8 bits and 4:2:0 chroma subsampling, obtained from the XIPH
repository \cite{xipha1}. Each sequence was encoded at four quantization points
(\texttt{cq-level} 20, 32, 43, and 55) over five frames sampled uniformly along
the whole clip, since consecutive frames are nearly identical under All-Intra. The
resulting dataset holds 26.98 million samples, of which 10.07 million are decision
nodes. The split is made by sequence and was frozen before any measurement, with
ten sequences for training, three for validation, and three for the reserved test
set. No video
sequence was reused across the training, testing, or evaluation datasets, thus
avoiding data leakage and overfitting. All results reported here are measured on
the six held-out sequences, which contribute 226\,447 super blocks and
3\,808\,703 decision nodes. The models were implemented in PyTorch \cite{pytorch}.

\subsection{Fair Treatment of the Deep Arm}

Since the deep arm carries the hypothesis that raw pixels plus capacity suffice, a
negative result on it is only meaningful if the arm was treated fairly, and three
measures were taken. First, it was retrained on the same corpus and under the same
split as the rungs, after an audit of an earlier checkpoint showed that it had
been trained on a smaller corpus with a single quantization point and with two of
the six held-out sequences contaminated. Second, capacity was controlled by
doubling the fusion width from 128 to 256 channels, and the wider model is 1.0\%
worse in validation loss, which establishes that capacity is not the binding
constraint.

Third, the training objective was ablated. A cost-sensitive cross entropy
weighting each node by $1 + \alpha\,\hat{r}(n)$, with $\hat{r}$ the normalized
loss of \eqref{eqn_regret}, was compared against plain cross entropy, obtained by
setting $\alpha = 0$. Under identical corpus, schedule, and selection criterion,
the cost-sensitive objective at $\alpha = 3$ is better at every operating point
from 10\% of cost reduction onward, and it is the deep result discussed in
Section~\ref{sec_results}. This inverts a conclusion previously drawn by the
authors from an uncontrolled comparison, so the value reported here is the most
favorable one measured for the deep arm.

\section{Results and Discussion}
\label{sec_results}

The representations were evaluated with the protocol of Section III over the six
held-out sequences and 3\,808\,703 decision nodes, under the All-Intra
configuration of libaom v3.10.0 at \texttt{cpu-used=0}. Table~\ref{tab_frontier}
reports the loss $L$ of \eqref{eqn_loss} at five matched cost reductions, in units
of $10^{-3}$ percent of the aggregated undivided cost, where lower is better, and
Fig.~\ref{fig_frontier} shows the corresponding frontiers.

\begin{table}[!t]
\renewcommand{\arraystretch}{1.3}
\caption{Rate-Distortion Loss at Matched Cost Reduction, in Units of
\upshape{$10^{-3}$} Percent of the Aggregated Undivided Cost}
\label{tab_frontier}
\centering
\begin{tabular}{lrrrrr}
\hline
\textbf{Representation} & \textbf{10\%} & \textbf{15\%} & \textbf{20\%}
& \textbf{25\%} & \textbf{30\%} \\
\hline
Random control & 196.3 & 298.9 & 406.6 & 516.6 & 634.2 \\
Variance & 4.3 & 12.7 & 25.0 & 37.4 & 55.5 \\
ConvNeXt, plain CE & 6.5 & 10.5 & 16.4 & 27.2 & 44.2 \\
ConvNeXt, width 256 & 7.5 & 12.2 & 19.4 & 29.4 & 44.4 \\
ConvNeXt, cost-sensitive & 5.4 & 8.0 & 14.0 & 21.9 & 37.9 \\
A (24 columns) & 0.5 & 1.3 & 2.8 & 5.3 & 9.1 \\
A+C (28 columns) & 0.5 & 1.2 & 2.7 & 6.4 & 12.2 \\
A+B+C (36 columns) & 0.2 & 0.6 & 1.8 & 4.0 & 7.8 \\
A+B (32 columns) & \textbf{0.2} & \textbf{0.5} & \textbf{1.5} & \textbf{3.3}
& \textbf{6.9} \\
\hline
\end{tabular}
\end{table}

\begin{figure}[!t]
\centering
% TODO: replace by fig2_frontier.pdf, generated last in av1_bench.
\framebox[0.98\columnwidth]{\rule{0pt}{5.4cm}%
\parbox{0.86\columnwidth}{\centering\footnotesize
PLACEHOLDER --- rate-distortion loss versus matched cost reduction, logarithmic
vertical axis, one curve per representation, with seed-amplitude bars on the four
rungs.}}
\caption{Rate-distortion loss as a function of matched search-cost reduction. Bars
on the rungs show the amplitude over three training seeds.}
\label{fig_frontier}
\end{figure}

The main conclusion when observing Table~\ref{tab_frontier} is that the ordering
of the representations is preserved across the whole range, so it is not an
artifact of one operating point. At 25\%, isolated variance loses 37.4
units while the twenty-four compact columns of rung A lose 5.3, a reduction of 7.0
times, against the 516.6 units of the random control. This is the largest single
gain in the table, and it contradicts any reading in which the information
available in the pixel domain would be exhausted by variance.

Adding the eight causal partitioning columns to rung A lowers the loss from 5.3 to
3.3 units, a further reduction of 1.60 times. The separation is robust to
initialization, since the worst of the three seeds of A+B, at 3.6, is still better
than the best of the three seeds of A, at 4.9. This is the central result of this
work, because those eight columns describe the partition choices of the already
coded neighbors and are not computable from the luminance of the block under any
transformation. The information they carry is therefore produced by the encoding
process itself.

Conversely, the four columns of effective quantization, frame position, and node
depth add nothing. Rung A+C reaches 6.4 units against 5.3 for rung A at 25\%, and
its seed amplitude is 4.5 units against 0.8 for rung A, so the two distributions
overlap and the statement supported by the data is that block C does not add
information while destabilizing the fit. Once combined with the causal columns the
effect is no longer ambiguous, as A+B reaches 3.3 units against 4.0 for A+B+C, all
nine pairwise seed comparisons agree, and Fig.~\ref{fig_frontier} shows A+B
dominating A+B+C across the entire frontier. The thirty-six column vector is therefore not the offline optimum,
although this superiority of A+B has not been verified inside the encoder.

The deep arm does not support the hypothesis that raw pixels plus representation
capacity suffice. In its most favorable configuration it reaches 21.9 units at
25\%, which is 4.1 times worse than rung A, although it holds 28\,128\,638
parameters against the 11\,337 of the three networks of rung A, a factor of 2481.
Doubling the fusion width degrades the result rather than improving it, at 29.4
against 27.2 units under the same objective. This may be attributed to the
information that is missing rather than to the capacity that is available, since
the deep arm reads the same pixels as rung A and the gap is not closed by
capacity, whereas eight columns of encoder state close a further factor of 1.60
with 4291 parameters per level. It is important to emphasize that the deep arm is
therefore not an upper bound of the pixel domain, since a model outperformed by
another with strictly smaller information access establishes no performance
ceiling. The only genuine upper bound here is the reference search itself, whose
loss is zero by construction.

Three limitations qualify these results and are reported as part of them. The
protocol is ordinal, as established in Section III-B, so none of these values
converts into a Bj\o ntegaard delta bitrate or predicts what an encoder would
deliver. The three seeds provide an amplitude and not a confidence interval. The
objective ablation of the deep arm covers two points, $\alpha = 0$ and
$\alpha = 3$, and the curve in $\alpha$ was not swept. The comparison is
nevertheless conservative in the direction that matters, since every one of these
choices was resolved in favor of the deep arm.

\section{Conclusions}

This paper presented RPP, an offline attribution protocol that measures the
rate-distortion loss of competing representations of the AV1 intra partition
decision at matched search-cost reduction, reaching the conclusion that the causal
state of the encoding process carries decision information that luminance alone
does not provide. The protocol holds the pruning policy and the training recipe
fixed and varies only the information a representation may read, charging every
forced decision the overhead recorded by an exhaustive reference search. It was
evaluated over 3.81 million held-out decision nodes of sixteen UVG 4K sequences
encoded with libaom v3.10.0. Twenty-four compact descriptors
improve on isolated variance by 7.0 times, eight causal partitioning columns
improve on those by a further 1.60 times, and four quantization and position
columns add nothing. A ConvNeXt with 28.1 million parameters reading raw pixels
remains 4.1 times behind the compact descriptors, and doubling its width makes it
worse, so capacity is not the binding constraint. Besides, to the best of the
authors' knowledge, this is the first work to attribute rate-distortion loss to
information access in AV1 partitioning.

\section*{Acknowledgment}

This study was financed in part by the Coordena\c{c}\~ao de Aperfei\c{c}oamento de
Pessoal de N\'ivel Superior -- Brasil (CAPES) -- Finance Code 001, and also by the
Brazilian research support agencies CNPq and FAPERGS.

\section*{Compliance with Ethical Standards}

This work is a computational study over publicly available video sequences and
does not involve human participants or animals. The authors declare that they have
no conflict of interest.

% Manual column equalization of the last page (IEEE requirement): forces the
% column break before this reference. Recheck if the reference list changes.
\IEEEtriggeratref{18}
\begin{thebibliography}{00}

\bibitem{statista} Statista. (2025). \emph{Share of Internet Users Watching Online Videos Every Month From 1st Quarter 2022 to 3rd Quarter 2024}. [Online]. Available: \url{https://www.statista.com/statistics/1489445/online-video-viewers-worldwide-quarterly/}

\bibitem{hevc} ITU-T. (2013). \emph{H.265: High Efficiency Video Coding}. [Online]. Available: \url{https://www.itu.int/rec/t-rec-h.265}

\bibitem{vvc} ITU-T. (2020). \emph{H.266: Versatile Video Coding}. [Online]. Available: \url{https://www.itu.int/rec/t-rec-h.266}

\bibitem{av1site} Alliance for Open Media. (2019). \emph{AV1 Video Codec}. [Online]. Available: \url{https://aomedia.org/av1}

\bibitem{han2021} J. Han \emph{et al.}, ``A technical overview of AV1,'' \emph{Proc. IEEE}, vol. 109, no. 9, pp. 1435--1462, Sep. 2021.

\bibitem{chen2018} Y. Chen \emph{et al.}, ``An overview of core coding tools in the AV1 video codec,'' in \emph{Proc. Picture Coding Symp. (PCS)}, Jun. 2018, pp. 41--45.

\bibitem{grois2016} D. Grois, T. Nguyen, and D. Marpe, ``Coding efficiency comparison of AV1/VP9, H.265/MPEG-HEVC, and H.264/MPEG-AVC encoders,'' in \emph{Proc. Picture Coding Symp. (PCS)}, Dec. 2016, pp. 1--5.

\bibitem{layek2017} M. A. Layek \emph{et al.}, ``Performance analysis of H.264, H.265, VP9 and AV1 video encoders,'' in \emph{Proc. 19th Asia-Pacific Netw. Oper. Manage. Symp. (APNOMS)}, Sep. 2017, pp. 322--325.

\bibitem{sullivan1998} G. J. Sullivan and T. Wiegand, ``Rate-distortion optimization for video compression,'' \emph{IEEE Signal Process. Mag.}, vol. 15, no. 6, pp. 74--90, Nov. 1998.

\bibitem{xu2020} M. Xu and B. Jeon, ``Selection of intra prediction tools for fast AV1 encoding,'' in \emph{Proc. IEEE Int. Symp. Broadband Multimedia Syst. Broadcast. (BMSB)}, Oct. 2020, pp. 1--4.

\bibitem{xu2021} M. Xu and B. Jeon, ``User-priority based AV1 coding tool selection,'' \emph{IEEE Trans. Broadcast.}, vol. 67, no. 3, pp. 736--745, Sep. 2021.

\bibitem{jeong2019} J. Jeong, G. Gankhuyag, and Y.-H. Kim, ``A fast intra mode decision based on accuracy of rate distortion model for AV1 intra encoding,'' in \emph{Proc. 34th Int. Tech. Conf. Circuits/Syst., Comput. Commun. (ITC-CSCC)}, Jun. 2019, pp. 1--3.

\bibitem{kim2019} J. Kim, S. Blasi, A. S. Dias, M. Mrak, and E. Izquierdo, ``Fast inter-prediction based on decision trees for AV1 encoding,'' in \emph{Proc. IEEE Int. Conf. Acoust., Speech Signal Process. (ICASSP)}, May 2019, pp. 1627--1631.

\bibitem{saldanha2022} M. Saldanha, G. Sanchez, C. Marcon, and L. Agostini, ``Configurable fast block partitioning for VVC intra coding using light gradient boosting machine,'' \emph{IEEE Trans. Circuits Syst. Video Technol.}, vol. 32, no. 6, pp. 3947--3960, Jun. 2022.

\bibitem{he2021} Q. He, W. Wu, L. Luo, C. Zhu, and H. Guo, ``Random forest based fast CU partition for VVC intra coding,'' in \emph{Proc. IEEE Int. Symp. Broadband Multimedia Syst. Broadcast. (BMSB)}, Aug. 2021, pp. 1--4.

\bibitem{chen2019} J. Chen, H. Sun, J. Katto, X. Zeng, and Y. Fan, ``Fast QTMT partition decision algorithm in VVC intra coding based on variance and gradient,'' in \emph{Proc. IEEE Vis. Commun. Image Process. (VCIP)}, Dec. 2019, pp. 1--4.

\bibitem{correa2020} M. Corr\^ea, B. Zatt, D. Palomino, G. Corr\^ea, and L. Agostini, ``A high-throughput hardware architecture for AV1 non-directional intra modes,'' \emph{IEEE Trans. Circuits Syst. I, Reg. Papers}, vol. 67, no. 5, pp. 1481--1494, May 2020.

\bibitem{netto2022} L. Netto, M. Corr\^ea, D. Palomino, L. Agostini, and G. Corr\^ea, ``Power-quality configurable hardware design for AV1 directional intra-frame prediction,'' \emph{IEEE Design Test}, vol. 39, no. 2, pp. 38--45, Apr. 2022.

\bibitem{convnext} Z. Liu, H. Mao, C.-Y. Wu, C. Feichtenhofer, T. Darrell, and S. Xie, ``A ConvNet for the 2020s,'' in \emph{Proc. IEEE/CVF Conf. Comput. Vis. Pattern Recognit. (CVPR)}, Jun. 2022, pp. 11976--11986.

\bibitem{adamw} I. Loshchilov and F. Hutter, ``Decoupled weight decay regularization,'' in \emph{Proc. Int. Conf. Learn. Represent. (ICLR)}, May 2019, pp. 1--18.

\bibitem{pytorch} A. Paszke \emph{et al.}, ``PyTorch: An imperative style, high-performance deep learning library,'' in \emph{Proc. Adv. Neural Inf. Process. Syst. (NeurIPS)}, Dec. 2019, pp. 8024--8035.

\bibitem{libaom} Alliance for Open Media. \emph{AV1 Codec Library}, v3.10.0. [Online]. Available: \url{https://aomedia.googlesource.com/aom}

\bibitem{ctc} X. Zhao \emph{et al.}, ``AOM common test conditions v9.0,'' Alliance for Open Media, Codec Working Group, doc. CWG-G082, Mar. 2026.

\bibitem{uvg} A. Mercat, M. Viitanen, and J. Vanne, ``UVG dataset: 50/120fps 4K sequences for video codec analysis and development,'' in \emph{Proc. 11th ACM Multimedia Syst. Conf. (MMSys)}, Jun. 2020, pp. 297--302.

\bibitem{xipha1} Xiph.Org Foundation. \emph{AOM CTC Test Set, Class A1 (4K)}. [Online]. Available: \url{https://media.xiph.org/video/aomctc/test_set/a1_4k/}

\bibitem{bjontegaard} G. Bj{\o}ntegaard, ``Calculation of average PSNR differences between RD-curves,'' Video Coding Experts Group (VCEG), doc. VCEG-M33, Mar. 2001.

\end{thebibliography}

\end{document}
